\documentclass[12pt]{article}
\usepackage{fontspec}
\usepackage{graphicx}
\usepackage{pst-plot}
\usepackage{scrextend}
\usepackage{advdate}
\usepackage{listings}
\usepackage{tikz}
\usepackage{amssymb}
\usepackage{color}
\usepackage{soul}
\usepackage{caption}
\usepackage{fancyhdr}
\usepackage{pgfornament}
\usepackage[many]{tcolorbox}
\usepackage{collectbox}
\usepackage{mdframed}
\usepackage{xcolor}
\usepackage{mathtools}
\usepackage[hidelinks]{hyperref}
\usepackage[final]{pdfpages}
\usepackage[left=0.5in,right=0.5in,top=0.5in,bottom=0.5in,footskip=15mm]{geometry}

\newfontfamily{\ubuntu}[Path=./fonts/,]{UbuntuMono-Regular}
\newfontfamily{\ubuntui}[Path=./fonts/,]{UbuntuMono-Italic}
\newfontfamily{\ubuntub}[Path=./fonts/,]{UbuntuMono-Bold}

\definecolor{blue1}{RGB}{115, 3, 252}
\definecolor{blue2}{RGB}{3, 53, 252}
\definecolor{orange1}{RGB}{252, 107, 3}

\begin{document}
\pagenumbering{gobble}

\vspace{5mm}

\begin{tcolorbox}[width=0.41\textwidth, colback={white}, coltext=black, colframe={orange1}, fontupper=\bfseries]
{\ubuntu \textbf{DUE:} Wednesday, April 12, 2023}
\end{tcolorbox} \hfill {\ubuntu\textbf{\color{blue1}Rocco DeVito}}

\hfill {\ubuntu Issued \AdvanceDate[0]\today}

\vspace{10mm}

\noindent {\Large \textbf{\textcolor{blue1}{{\ubuntu Modular Arithmetic and} $(\mathbb{Z}_{12}, +)$ \hrulefill}}}

\vspace{10mm}

\noindent {\ubuntu Imagine you live in Clock World, where you can only do addition using the hour hand of a clock.}

\vspace{5mm}

\begin{figure}[h] \centering
    \includegraphics[width = 0.8\textwidth]{images/clock_world.jpeg}
    \caption{\textcolor{blue1}{{\ubuntu Clock World: A.K.A. Salvador Dali's \ubuntui{The Persistence of Memory}}}}
    \label{fig:dali}
\end{figure}

\vspace{5mm}

\noindent {\ubuntu You would be limited in how many numbers you had to work with. To see why, imagine that the hour hand is currently pointing to the number $11$. Now advance by $15$ hours. Would the hour hand be pointing to the number $26$ (which is what $11 + 15$ is equal to)? What do you think? If you are unsure, look at the diagram below.}

\begin{figure}[h] \centering
    \includegraphics[width = 0.9\textwidth]{images/clock_addition_11_plus_15.png}
    \caption{\textcolor{blue1}{{\ubuntu On a clock, $11 + 15 = 2$}}}
    \label{fig:example1}
\end{figure}

\pagebreak

\noindent {\ubuntu The number $26$ isn't actually on a clock, so it can't be the number the hour hand is pointing to. Whenever the hour hand goes beyond $12$, the count starts over again. If this makes sense to you, try to solve the following Clock World addition questions.}

\vspace{5mm}

\begin{figure}[h]
    \centering
    \begin{minipage}{0.65\textwidth}
        \includegraphics[width=\textwidth]{images/clock_question_3_plus_17.png}
    \end{minipage}
    \hfill
\begin{minipage}{0.3\textwidth}
        \begin{tikzpicture}
            \draw[dashed] (0,0) rectangle (\textwidth, \textwidth); % Outlined square
        \end{tikzpicture}
    \end{minipage}
\end{figure}

\vspace{5mm}

\begin{figure}[h]
    \centering
    \begin{minipage}{0.65\textwidth}
        \includegraphics[width=\textwidth]{images/clock_question_5_plus_12.png}
    \end{minipage}
    \hfill
\begin{minipage}{0.3\textwidth}
        \begin{tikzpicture}
            \draw[dashed] (0,0) rectangle (\textwidth, \textwidth); % Outlined square
        \end{tikzpicture}
    \end{minipage}
\end{figure}

\vspace{5mm}

\begin{figure}[h]
    \centering
    \begin{minipage}{0.65\textwidth}
        \includegraphics[width=\textwidth]{images/clock_question_8_plus_26.png}
    \end{minipage}
    \hfill
\begin{minipage}{0.3\textwidth}
        \begin{tikzpicture}
            \draw[dashed] (0,0) rectangle (\textwidth, \textwidth); % Outlined square
        \end{tikzpicture}
    \end{minipage}
\end{figure}

\pagebreak

\noindent {\ubuntu You might think that using pictures of clocks is a cumbersome way of doing mathematics. I agree, and so does every sane mathematician on earth. That's why we have mathematical notation. Allow me to introduce you to the mathematical object $(\mathbb{Z}_{12}, +).$ This mathematical object is technically a} \ubuntub{group}, \ubuntu{but we will get into what groups are at a later point. Let's examine the components of this object.}

\vspace{5mm}

\noindent {\ubuntu The notation $\mathbb{Z}_{12}$ represents the set of 'integers' from $0$ to $11$ (technically, they are \textbf{equivalence classes}, not integers, but don't worry about that yet). To make things clear for you, you can think of $\mathbb{Z}_{12}$ as}

\vspace{3mm}

\begin{center}
\scalebox{1.5}{$\mathbb{Z}_{12} = \left\lbrace 0,1,2,3,4,5,6,7,8,9,10,11\right\rbrace.$}
\end{center}

\vspace{3mm}

\noindent {\ubuntu The only difference between the numbers in $\mathbb{Z}_{12}$ and the numbers on a clock is that $\mathbb{Z}_{12}$ has a 0 where a clock has $12$. But, in Clock World, $12$ and $0$ are equal, so this isn't a huge problem. Here's a new clock, just to make this a little clearer.}

\begin{figure}[h] \centering
    \includegraphics[width = 0.4\textwidth]{images/mod12_clock.png}
    \caption{\textcolor{blue1}{{\ubuntu In Clock World, $12 = 0$}}}
    \label{fig:example1}
\end{figure}

\noindent {\ubuntu If you are familiar with military time, then you know that 12AM is equivalent to 00:00. This might help your understanding. You might also know that 1PM is equivalent to 13:00 and 2PM is equivalent to 14:00. Likewise, in $\mathbb{Z}_{12}$, since we only have the integers $0$ to $11$ to work with, $13 = 1$ and $14 = 2$. What would $17$ be equal to?}

\vspace{5mm}

\noindent {\ubuntu As for the $+$ in $(\mathbb{Z}_{12}, +)$, it represents something called 'addition modulo $12$.' You can think of it as 'wrap-around' addition, i.e., addition such that you can never get a sum higher than $11$. If you added two numbers together and got $12$, you would wrap-around to $0$; if you added two numbers together and got $13$, you would wrap-around to $1$; etc.}

\vspace{5mm}

\noindent {\ubuntu So the mathematical object $(\mathbb{Z}_{12}, +)$ is a model of a clock whose hour hand can be moved discretely in a clockwise direction. Here, 'discretely' means that the hour hand can't land on $8.5$ or $11.5678$; it can only land on integers. In other words, we aren't moving the hour hand through a continuum, if this word means anything to you. If it doesn't, don't worry about it.}

\pagebreak

\noindent {\ubuntu If I wanted to ask you to perform a Clock World addition question without using clock pictures, I would write something like this:}

\begin{center}
    \scalebox{1.5}{{\ubuntu Find the sum $4 + 25 \text{ (mod } 12).$}}
\end{center}

\noindent {\ubuntu This question is equivalent to the following picture question:}

\begin{figure}[h]
    \centering
    \begin{minipage}{0.65\textwidth}
        \includegraphics[width=\textwidth]{images/notation_vs_pic.png}
    \end{minipage}
    \hfill
\begin{minipage}{0.3\textwidth}
        \begin{tikzpicture}
            \draw[dashed] (0,0) rectangle (\textwidth, \textwidth); % Outlined square
        \end{tikzpicture}
    \end{minipage}
\end{figure}

\vspace{10mm}

\hrulefill

\vspace{10mm}

\begin{center}
    {\Large \textbf{\textcolor{blue1}{Exercises}}}
\end{center}

\noindent{\ubuntu\textbf{Instructions:} Compute the following sums modulo $12$.}

\vspace{15mm}

\noindent \scalebox{2}{\textcolor{blue1}{(1.)} $2 + 13 \text{ (mod } 12) = $}

\vspace{12mm}

\noindent \scalebox{2}{\textcolor{blue1}{(2.)} $9 + 24 \text{ (mod } 12) = $}

\vspace{12mm}

\noindent \scalebox{2}{\textcolor{blue1}{(3.)} $0 + 36 \text{ (mod } 12) = $}

\vspace{12mm}

\noindent \scalebox{2}{\textcolor{blue1}{(4.)} $7 + 19 \text{ (mod } 12) = $}

\vspace{12mm}

\noindent \scalebox{2}{\textcolor{blue1}{(5.)} $1 + 27 \text{ (mod } 12) = $}


\end{document}
